\section{Related Work}
\label{sec::related}

\subsection{Classical Clustering Methods for scRNA-seq}
Numerous single-cell clustering methods have been proposed in recent years. Early approaches typically follow a two-stage paradigm: first obtaining low-dimensional features via dimensionality reduction techniques, and subsequently applying classical algorithms for clustering, such as  $k$-means~\cite{hicks2021mbkmeans}, hierarchical clustering~\cite{johnson1967hierarchical}, and density-based methods~\cite{petegrosso2020machine}. To learn low-dimensional representations, some methods, such as pcaReduce~\cite{vzurauskiene2016pcareduce}, integrate Principal Component Analysis (PCA) with $k$-means to iteratively merge clusters. In contrast, CIDR~\cite{lin2017cidr} reduces dimensionality by constructing a dissimilarity matrix based on imputed gene expression values. Nonetheless, these methods lack the flexibility to model complex nonlinear relationships among cells~\cite{stegle2015computational}. To overcome the limitations, various deep clustering methods have been developed to learn more expressive representations. DEC~\cite{xie-et-al:dec} pioneers this direction by learning feature representations and cluster assignments within a unified framework. Furthermore, scDeepCluster~\cite{tian-et-al:model-based-scrnaseq} and scDCC~\cite{tian2021model} incorporate Zero-Inflated Negative Binomial (ZINB) autoencoders, whereas contrastive-sc~\cite{ciortan2021contrastive} leverages contrastive learning for robustness. Despite their success, these methods primarily focus on intra-cellular gene expression patterns and often neglect topological relationships among cells.


\subsection{Deep Graph Clustering for scRNA-seq}
Recent research has pivoted toward deep graph clustering frameworks based on Graph Neural Networks (GNNs)~\cite{ju2023glcc,wang2025deep,ren2025mhgc,zhang2026fairgc,ju2026compactness}. By representing cells as nodes and their interactions as edges, these methods effectively integrate gene expression profiles with topological structures via message-passing mechanisms to learn context-aware cell embeddings. For instance, scGNN~\cite{wang2021scgnn} employs GNNs and multi-modal autoencoders to aggregate cell-cell relationships, while scGAE~\cite{luo2021topology} utilizes a multitask-oriented graph autoencoder to simultaneously preserve feature and topological information. To better handle the prevalence of dropout events and technical noise, scTAG~\cite{yu2022zinb} integrates a ZINB model into a topology-adaptive graph convolutional autoencoder to learn low-dimensional latent representations in an end-to-end manner. Similarly, scDSC~\cite{gan2022deep} combines a ZINB-based autoencoder with GNN modules through a mutually supervised strategy, allowing the structural and attribute information to mutually refine the clustering results. Most recently, scSiamese~\cite{xu2025scsiameseclu} leverages a Siamese clustering framework to enhance representation robustness via cross-view consistency. However, these GNN-based approaches primarily focus on local neighborhood aggregation and often overlook richer structural information, which results in over-smoothed embeddings that fail to distinguish biologically distinct populations~\cite{li2018deeper}.

